% results_log.tex — running log of quantization PPL results
% Evaluation: WikiText-2 perplexity, W4A4 (true weight + activation quantization)
%   Weights:     NVFP4 (E2M1, per-block absmax/6 scale, Hadamard-rotated)
%   Activations: Hadamard-rotated, then quantized per scheme (see table notes)
% Windows: 10 × 2048 tokens

\documentclass{article}
\usepackage{booktabs}
\usepackage{siunitx}
\usepackage{xcolor}
\usepackage{hyperref}
\usepackage[margin=1in]{geometry}

\sisetup{
  table-number-alignment = right,
  table-format = 2.4,
}

\definecolor{good}{HTML}{1A9E94}
\definecolor{moderate}{HTML}{A86A18}
\definecolor{critical}{HTML}{B83535}

\title{NVFP4 W4A4 Quantization --- PPL Results Log}
\author{Jay Vap}
\date{Updated: \today}

\begin{document}
\maketitle

% ─────────────────────────────────────────────────────────────────
\section{OPT-1.3B (WikiText-2 PPL)}
% ─────────────────────────────────────────────────────────────────

FP16 baseline: \textbf{14.8507}

\begin{table}[h]
\centering
\caption{OPT-1.3B --- W4A16 weight quantization, WikiText-2 PPL}
\label{tab:opt1.3b}
\begin{tabular}{l S[table-format=2.4] S[table-format=+2.4] S[table-format=+3.1,table-sign-mantissa]}
\toprule
Scheme & {PPL} & {$\Delta$PPL} & {Rel.\,overhead (\%)} \\
\midrule
fp16\_baseline            & 14.8507 & {---}   & {---}   \\
\addlinespace[2pt]
gf4                       & 17.3563 & +2.5055 & +16.9   \\
nvfp4\_absmax             & 18.6329 & +3.7822 & +25.4   \\
nvfp4\_norm$^\dagger$     & 22.2507 & +7.4000 & +49.8   \\
e2m1\_rms\_best           & 22.3748 & +7.5240 & +50.6   \\
nvfp4\_raw$^\ddagger$     & 59.9681 & +45.117 & +303.8  \\
\bottomrule
\end{tabular}
\end{table}

$^\dagger$ Bug\,1 fix: E2M1 magnitudes normalised by 6.0 before codebook lookup.\\
$^\ddagger$ Bug\,1: raw E2M1 magnitudes $\{0, 0.5, \ldots, 6.0\}$ passed to $[0,1]$-clamped
quantiser collapse 8 levels to 3 effective bins.

% ─────────────────────────────────────────────────────────────────
\section{OPT-13B (WikiText-2 PPL)}
% ─────────────────────────────────────────────────────────────────

FP16 baseline: \textbf{10.0913}

\begin{table}[h]
\centering
\caption{OPT-13B --- W4A4 (weights + activations), WikiText-2 PPL}
\label{tab:opt13b}
\begin{tabular}{l S[table-format=2.4] S[table-format=+2.4] S[table-format=+3.1,table-sign-mantissa]}
\toprule
Scheme & {PPL} & {$\Delta$PPL} & {Rel.\,overhead (\%)} \\
\midrule
fp16\_baseline            & 10.0913 & {---}   & {---}   \\
\addlinespace[2pt]
gf4                       & 10.8138 & +0.7225 & +7.1    \\
nvfp4\_w4a4$^\star$       & 10.9229 & +0.8317 & +8.2    \\
nvfp4\_absmax             & 11.0106 & +0.9194 & +9.1    \\
nvfp4\_norm$^\dagger$     & 11.9221 & +1.8308 & +18.1   \\
e2m1\_rms\_best           & 11.9075 & +1.8162 & +18.0   \\
nvfp4\_raw$^\ddagger$     & 44.9463 & +34.855 & +345.4  \\
\bottomrule
\end{tabular}
\end{table}

$^\star$ Canonical NVFP4 W4A4: block-16 E4M3 scale for both weights and activations
(Hadamard-rotated). Matches NVIDIA Blackwell FP4 hardware spec.
Only 0.11\,PPL behind GF4 on OPT-13B. OPT-1.3B pending.

% ─────────────────────────────────────────────────────────────────
\section{Cross-Scale Summary}
% ─────────────────────────────────────────────────────────────────

\begin{table}[h]
\centering
\caption{Relative PPL overhead (\%) by model size --- smaller is better}
\label{tab:crossscale}
\begin{tabular}{l r r r}
\toprule
Scheme & OPT-1.3B & OPT-13B & Scale factor \\
\midrule
gf4                   & +16.9\%  & +7.1\%   & 0.42$\times$ \\
nvfp4\_w4a4$^\star$   & pending  & +8.2\%   & —            \\
nvfp4\_absmax         & +25.4\%  & +9.1\%   & 0.36$\times$ \\
nvfp4\_norm           & +49.8\%  & +18.1\%  & 0.36$\times$ \\
e2m1\_rms\_best       & +50.6\%  & +18.0\%  & 0.36$\times$ \\
nvfp4\_raw            & +303.8\% & +345.4\% & 1.14$\times$ \\
\bottomrule
\end{tabular}
\end{table}

Scale factor = (13B rel.\,overhead) / (1.3B rel.\,overhead).
All working schemes improve as model size increases.
\texttt{nvfp4\_w4a4} (+8.2\% on 13B) is only 1.1\,pp behind GF4,
confirming that block-16 E4M3 scaling is the key to competitive NVFP4 W4A4.
$^\star$ OPT-1.3B run pending.

% ─────────────────────────────────────────────────────────────────
\section{Scheme Descriptions}
% ─────────────────────────────────────────────────────────────────

\begin{description}
  \item[\texttt{gf4}] W4(NVFP4) + A4(GF4). Weights in E2M1/Hadamard; activations
    quantized with Lloyd-Max Gaussian codebook + RMS$\times2.5$ scale (working baseline).
  \item[\texttt{nvfp4\_raw}] W4(NVFP4) + A4(E2M1\,raw). Both weights and activations
    in E2M1, but activation quantizer receives raw magnitudes $\{0,\ldots,6\}$ into
    a $[0,1]$-clamped lookup --- only 3 of 8 activation codebook levels reachable.
    \textbf{Bug\,1 in activations}.
  \item[\texttt{nvfp4\_norm}] W4(NVFP4) + A4(E2M1/6). Bug\,1 fix: activation magnitudes
    normalised by $\div6$ before lookup; RMS$\times2.5$ per-block scale.
    True symmetric NVFP4 W4A4 with RMS activation scaling.
  \item[\texttt{nvfp4\_absmax}] W4(NVFP4) + A4(E2M1/6, absmax). Same as above but
    activations use per-block absmax/$\max(\text{codebook})$ scale.
    \textbf{Canonical NVFP4 W4A4 baseline} (absmax activation scale, no RMS).
    Surprisingly competitive on 13B.
  \item[\texttt{e2m1\_rms\_best}] W4(NVFP4) + A4(E2M1, bias-searched).
    Exponent-bias grid search over $\{0,1,2\}$ per activation block; best-MSE bias kept.
    Best-effort NVFP4 W4A4 reference.
\end{description}

% ─────────────────────────────────────────────────────────────────
\section{Known Bugs}
% ─────────────────────────────────────────────────────────────────

\paragraph{Bug 1 (3-level collapse).}
Raw E2M1 magnitudes range $[0, 6]$.  The GF4 quantiser clips inputs to $[0,1]$, so only
thresholds 0.25 and 0.75 are reachable.  This leaves three bins: $\{0\text{--}0.25,\;
0.25\text{--}0.75,\;0.75\text{--}1.0\}$, effectively collapsing 8 codebook levels to 3.
Fix: normalise by $\div 6$ before lookup.

\paragraph{Bug 2 (padding absmax).}
On OPT-13B with padded sequences, extreme activation values on padding tokens dominate
the per-block absmax, compressing real-token activations to low codebook levels.
\texttt{nvfp4\_absmax} is likely mitigated at 13B scale by the model's greater robustness.

\end{document}
